% 链式法则（综述）
% license CCBYSA3
% type Wiki

本文根据 CC-BY-SA 协议转载翻译自维基百科\href{https://en.wikipedia.org/wiki/Chain_rule}{相关文章}

在微积分中，链式法则是一条公式，它用函数$f$和$g$的导数来表达它们复合函数的导数。更准确地说，若$h = f \circ g$是函数，满足$h(x) = f(g(x))$对任意 $x$ 成立，则在拉格朗日记号中，链式法则为：
$$
h'(x) = f'(g(x)) \, g'(x).~
$$
或者等价地：
$$
h' = (f \circ g)' = (f' \circ g)\cdot g'.~
$$
链式法则也可以用莱布尼茨记号表示。若变量 $z$ 依赖于变量 $y$，而 $y$ 又依赖于变量 $x$（即 $y$ 和 $z$ 都是依赖变量），那么 $z$ 也通过中间变量 $y$ 依赖于 $x$。在这种情况下，链式法则表示为：
$$
\frac{dz}{dx} = \frac{dz}{dy} \cdot \frac{dy}{dx},~
$$
以及：
$$
\left.\frac{dz}{dx}\right|_{x} = \left.\frac{dz}{dy}\right|_{y(x)} \cdot \left.\frac{dy}{dx}\right|_{x},~
$$
其中竖线记号用来指出导数应当在哪些点进行计算。

在积分中，与链式法则对应的规则是代换积分法。
\subsection{直观解释}
直观上，链式法则表明：如果已知 $z$ 相对于 $y$ 的瞬时变化率，以及 $y$ 相对于 $x$ 的瞬时变化率，那么就可以通过两者的乘积来计算 $z$ 相对于 $x$ 的瞬时变化率。

正如 George F. Simmons 所说：“如果一辆汽车的速度是自行车的两倍，而自行车的速度是步行者的四倍，那么汽车的速度就是步行者的 $2 \times 4 = 8$ 倍。”【1】

这个例子与链式法则的关系如下：设 $z, y, x$ 分别表示汽车、自行车和步行者的位置（变量）。汽车与自行车相对位置的变化率为：$\frac{dz}{dy} = 2$.同样，自行车与步行者相对位置的变化率为：$\frac{dy}{dx} = 4$.因此，汽车与步行者相对位置的变化率为：
$$
\frac{dz}{dx} = \frac{dz}{dy}\cdot \frac{dy}{dx} = 2 \cdot 4 = 8.~
$$
位置变化率就是速度之比，而速度就是位置对时间的导数。即：
$$
\frac{dz}{dx} = \frac{\tfrac{dz}{dt}}{\tfrac{dx}{dt}},~
$$
或者等价地：
$$
\frac{dz}{dt} = \frac{dz}{dx}\cdot \frac{dx}{dt},~
$$
这同样是链式法则的一个应用。
\subsection{历史}
链式法则最早似乎是由戈特弗里德·威廉·莱布尼茨使用的。他利用它来计算$\sqrt{a+bz+cz^{2}}$的导数，把它看作平方根函数与函数$a+bz+cz^{2}$的复合。他在 1676 年的一篇论文中首次提及这一点（其中计算中有一个符号错误）【2】。链式法则的常见记号归功于莱布尼茨【3】。纪尧姆·德·洛必达在其《无穷小分析》中隐含地使用了链式法则。尽管莱昂哈德·欧拉的分析著作写于莱布尼茨发现一百多年之后，但其中并未出现链式法则【需要引证】。一般认为，链式法则的第一个“现代”版本出现在拉格朗日 1797 年的《解析函数论》中；它也出现在柯西 1823 年的《在皇家综合理工学院所作无穷小计算讲义摘要》中【3】。
